\section{Literal mixed coefficient orders in the same relative chart}

Let $B=\mathcal O_{\mathbb C^d,0}$, $R=\mathcal O_{\mathbb C^{d+t},0}$,
and denote their genuine maximal ideals by $\mathfrak m_B,\mathfrak m_R$.
All normal coefficients below are the actual analytic germs extracted from
whole original normal-power quotients, not a chosen formal coefficient model.
The existing formal mixed-power theorem already counts tangent and normal
weights. The next results retain the actual analytic coefficient germs and
compute their literal restrictions in the SAME actual normalized chart.

\begin{theorem}[A whole actual point order is exactly its finite mixed analytic coefficient orders]
\label{hir:actual-mixed-analytic-coefficient-point-orders}
\uses{hir:actual-relative-coefficient-literal-formal-taylor-dictionary,hir:actual-original-formal-taylor-maximal-power-kernel}
\lean{TranscendentalNakamaye.Pending.HironakaActualMixedCoefficientOrders.mem_maximal_pow_iff_actual_normal_coefficients}
\leanok
For every original $g\in R$ and actual finite $q$,
\[
 g\in\mathfrak m_R^q\quad\Longleftrightarrow\quad
 \forall\gamma,\ |\gamma|<q\Longrightarrow
 c_\gamma(g)\in\mathfrak m_B^{q-|\gamma|}.
\]
These are finitely many whole coefficient-germ conditions, including $q=0$
and empty tangent or normal blocks.
\end{theorem}
\begin{proof}
The whole original Taylor power criterion, transported by its genuine dimension
reassociation, says that actual maximal-power membership is exactly vanishing
of all actual Taylor coefficients below the given total degree. The actual
Taylor normal-coefficient dictionary identifies the tangent coefficient at
$\beta$ of $\widehat{c_\gamma(g)}$ with the coefficient at the literal joined
exponent $(\beta,\gamma)$ of $\widehat g$. Its actual total degree is
$|\beta|+|\gamma|$. Thus a whole order $q$ kills every tangent coefficient
with $|\beta|<q-|\gamma|$, proving the displayed base power by the genuine
base Taylor criterion. Conversely split every exponent below $q$ into its
actual tangent and normal blocks. The latter has degree below $q$, and the
former has degree below the corresponding displayed base-power bound.
The assumed whole base power kills that exact Taylor coefficient. The original
whole maximal-power criterion gives $g\in\mathfrak m_R^q$.
\end{proof}

\begin{theorem}[The actual total homogeneous part has its literal finite mixed coefficient bands]
\label{hir:actual-mixed-homogeneous-coefficient-bands}
\uses{hir:actual-formal-centre-power-relative-coefficients}
\lean{TranscendentalNakamaye.Pending.HironakaActualMixedCoefficientOrders.actualFormalNormalCoefficient_homogeneousComponent}
\leanok
For an entire formal series $F$, its actual normal coefficient at $\gamma$
in the total homogeneous component of degree $q$ is zero when $|\gamma|>q$,
and otherwise is the actual tangent homogeneous component of degree
$q-|\gamma|$ of $c_\gamma(F)$. In particular these normal labels form a
finite bounded family; tangent-coefficient degrees are retained rather than
silently replaced by zero.
\end{theorem}
\begin{proof}
Compare the actual coefficient at every tangent exponent $\beta$.
On the left, the literal homogeneous-component formula retains precisely the
coefficient of $F$ whose joined exponent has degree $q$.
The actual joined degree is $|\beta|+|\gamma|$. If $|\gamma|>q$ this is
impossible and every coefficient is zero. Otherwise the equality is precisely
$|\beta|=q-|\gamma|$, which is the actual coefficient condition in the
indicated tangent homogeneous component. Whole formal-series extensionality
proves the entire coefficient-series identity.
\end{proof}

\begin{theorem}[The same actual normalized quotient gives whole pivot-free coefficient identities and finite mixed base orders]
\label{hir:actual-relative-quotient-mixed-fibre-coefficient-orders}
\uses{hir:actual-mixed-analytic-coefficient-point-orders,hir:actual-normal-quotient-canonical-exceptional-restriction,hir:actual-relative-homogeneous-dehomogenization-injective}
\lean{TranscendentalNakamaye.Pending.HironakaActualMixedCoefficientOrders.relative_quotient_pivot_free_coefficient,TranscendentalNakamaye.Pending.HironakaActualMixedCoefficientOrders.relative_quotient_mixed_coefficient_order}
\leanok
Fix actual original $G\in N^w$ and its SAME actual quotient
$e^wH=\pi^*G$ in a prescribed relative normal-pivot chart. Let
$p=\delta(f_w(G))\in B[z]$ be the literal dehomogenization of its actual
normal homogeneous polynomial. For every pivot-free exponent $\gamma$,
\[
 c_\gamma(H)=p_\gamma\qquad\text{as whole original analytic germs in }B.
\]
Consequently a natural CURRENT order $H\in\mathfrak m_R^r$ internally gives
$p_\gamma\in\mathfrak m_B^{r-|\gamma|}$ for every pivot-free
$|\gamma|<r$. No desired coefficient equality or mixed base-order estimate is
assumed, and no assertion of total-degree homogeneity follows from normal
homogeneity alone.
\end{theorem}
\begin{proof}
The genuine actual quotient restriction gives the WHOLE analytic congruence
$H-\operatorname{ev}(p)\in(e)$ for this exact quotient and this exact source
normal homogeneous polynomial. Its actual Taylor difference is divisible by
the genuine formal pivot variable. At every joined exponent $(\beta,\gamma)$
with zero pivot component, the actual variable-divisibility coefficient
criterion therefore makes that coefficient zero. The literal actual normal
coefficient dictionary for $H$ and for whole relative polynomial evaluation
identifies this with equality of every Taylor coefficient of the actual base
germs $c_\gamma(H)$ and $p_\gamma$. Actual base Taylor injectivity yields the
entire germ equality. Under the disclosed natural current-order hypothesis,
the first theorem produces the finite base power for $c_\gamma(H)$; substituting
the internally proved equality gives the claimed actual coefficient bound.
This is a concrete finite mixed coefficient constraint. Complete mixed-filtered
syzygy lifting, current maximal-symbol regularity and relative forward HS remain
subsequent obligations.
\end{proof}
